\documentclass[12pt]{article}
\usepackage{amsmath,amssymb,amsfonts,amsthm}
\usepackage{hyperref}
\usepackage{geometry}
\usepackage{authblk}
\geometry{margin=1in}

\newtheorem{theorem}{Theorem}[section]
\newtheorem{lemma}[theorem]{Lemma}
\newtheorem{proposition}[theorem]{Proposition}
\newtheorem{corollary}[theorem]{Corollary}
\theoremstyle{definition}
\newtheorem{definition}[theorem]{Definition}
\theoremstyle{remark}
\newtheorem{remark}[theorem]{Remark}

\newcommand{\R}{\mathbb{R}}
\newcommand{\Z}{\mathbb{Z}}
\newcommand{\Q}{\mathbb{Q}}
\newcommand{\End}{\mathrm{End}}
\newcommand{\Endzero}{\mathrm{End}^{\,\mathrm{tr}=0}}

\title{Discrete Shape Space and Integer Logarithmic Lattices}
uthor[1]{Marvin L.\ Gentry}
ffil[1]{Independent Researcher\
  	exttt{drmlgentry@protonmail.com}\
  ORCID: 0009-0006-4550-2663}
\date{}

\begin{document}
\maketitle

\begin{abstract}
Many physical and geometric observables are inherently multiplicative:
their ratios, not their absolute values, carry intrinsic meaning. When
expressed in logarithmic coordinates, an overall additive shift is
physically redundant. Removing this redundancy defines a quotient space
we call \emph{shape space}. We show that for quadratic descriptions of
rank~$n$, the number of independent scale-free deformation directions
is exactly $n^2-1$, the dimension of the traceless endomorphism space
$\mathrm{End}^{\,\mathrm{tr}=0}(V)$. For $n=3,4,5$ this yields $8,15,24$.
We establish this as a formal proposition with proof, note its coincidence
with adjoint representation dimensions of $\mathrm{SU}(n)$, and discuss
how golden-ratio-saturated geometries provide examples in which quadratic
invariants remain within the rank-2 integer module $\Z[\varphi]$.
\end{abstract}

\section{Introduction}

Many quantities of physical and geometric interest span many orders of
magnitude. Their ratios, not their absolute values, carry the meaningful
information. This suggests working in logarithmic coordinates, where
multiplicative relationships become additive. The space of logarithmic
coordinates carries a redundancy: an overall additive constant corresponds
to a simultaneous rescaling of all quantities. Factoring out this redundancy
leads to \emph{shape space}.

In this paper we study the structure of shape space for quadratic data.
Our central result (Proposition~\ref{prop:dim}) is that the space of
independent scale-free deformation directions for a rank-$n$ quadratic
description has dimension $n^2-1$, following from the rank-nullity theorem
applied to the trace map on endomorphisms.

\section{Logarithmic Coordinates and Shape Space}

\begin{definition}[Shape space]\label{def:shapespace}
The \emph{shape space} associated with $N$ positive quantities is the
quotient $\mathcal{S}_N=\R^N/\R\mathbf{1}$, where
$\R\mathbf{1}=\{(t,\ldots,t):t\in\R\}$. It has dimension $N-1$.
\end{definition}

\section{The $n^2-1$ Count}

Infinitesimal deformations of quadratic data on an $n$-dimensional space
are represented by endomorphisms $T\colon V\to V$. Uniform rescalings
$T=\lambda I$ are shape-redundant and must be quotiented out.

\begin{proposition}[Dimension of shape-space deformations]
\label{prop:dim}
The space of independent scale-free deformation directions of a rank-$n$
quadratic description is
\[
\dim\Endzero(V)=n^2-1,
\]
where $\Endzero(V)=\{T\in\End(V):\operatorname{tr}T=0\}$.
\end{proposition}

\begin{proof}
The trace map $\operatorname{tr}\colon\End(V)\to k$ is surjective with
kernel $\Endzero(V)$. By the rank-nullity theorem,
$\dim\Endzero(V)=\dim\End(V)-1=n^2-1$.
The canonical splitting $\End(V)=\Endzero(V)\oplus k\cdot I$ confirms
that $\Endzero(V)$ is exactly the scale-free quotient.
\end{proof}

\begin{corollary}
For ranks $n=3,4,5$ the independent deformation count is $8,15,24$.
\end{corollary}

\begin{remark}[Relation to $\mathrm{SU}(n)$ adjoint representations]
\label{rem:sun}
The integers $n^2-1$ are also the dimensions of the adjoint representations
of $\mathrm{SU}(n)$: $\dim\mathfrak{su}(3)=8$, $\dim\mathfrak{su}(4)=15$,
$\dim\mathfrak{su}(5)=24$. The Lie algebra $\mathfrak{su}(n)$ consists of
skew-Hermitian traceless $n\times n$ matrices, the complex analogue of the
real traceless endomorphisms here. The present derivation reaches the same
count from scale-free shape space without invoking any symmetry assumption~\cite{ratcliffe,weyl}.
\end{remark}

\section{Golden-Ratio Closure and Quadratic Invariants}

Many discrete geometries involve the golden ratio~\cite{coxeter} $\varphi=(1+\sqrt5)/2$,
which satisfies $\varphi^2=\varphi+1$. The set
$\Z[\varphi]=\{a+b\varphi:a,b\in\Z\}$ is a commutative ring (the ring
of integers of $\Q(\sqrt5)$~\cite{neukirch}).

\begin{proposition}[Quadratic closure of \texorpdfstring{$\Z[\varphi]$}{Z[phi]}]
\label{prop:closure}
Let $T\in M_n(\Z[\varphi])$. Then $\operatorname{tr}(T^2)\in\Z[\varphi]$.
More generally, all polynomial invariants of $T$ with integer coefficients
lie in $\Z[\varphi]$.
\end{proposition}

\begin{proof}
$\Z[\varphi]$ is a commutative ring closed under addition and multiplication.
The $(i,j)$ entry of $T^2$ is $\sum_k T_{ik}T_{kj}\in\Z[\varphi]$, and
$\operatorname{tr}(T^2)=\sum_i(T^2)_{ii}\in\Z[\varphi]$. The argument
for general polynomial invariants is identical.
\end{proof}

\begin{remark}[Scope of Proposition~\ref{prop:closure}]
Proposition~\ref{prop:closure} establishes that if deformation data lies
in $\Z[\varphi]$, then quadratic invariants also lie in $\Z[\varphi]$.
It does not select the values $n=3,4,5$; those come from
Proposition~\ref{prop:dim}. The role of $\Z[\varphi]$ is algebraic
closure, not geometric selection.
\end{remark}

\section{Orientation and Signed Directions}

\begin{definition}[Orientation-reversed deformation]
Let $T\in\Endzero(V)$. The \emph{orientation-reversed} deformation is
$-T\in\Endzero(V)$. In a Hermitian setting one may use the adjoint $T^\dagger$,
which reverses orientation while preserving norm.
\end{definition}

Signed integer coefficients in logarithmic coordinates have direct geometric
meaning: $+k$ represents $k$ steps forward along a deformation axis,
$-k$ represents $k$ steps in reverse.

\section{Scope and Generalization}

The derivation in Section~3 uses only finite dimensionality of $\End(V)$
and surjectivity of the trace map. It does not depend on any specific group,
representation, or geometric realization. For other discrete geometries
the role of $\Z[\varphi]$ in Proposition~\ref{prop:closure} is replaced by
the appropriate ring of integers of the relevant quadratic field; the
formula $n^2-1$ remains universal.

\section{Conclusion}

We have shown that $n^2-1$ arises as the dimension of the space of
independent scale-free deformation directions for a rank-$n$ quadratic
description, proved by rank-nullity alone. For $n=3,4,5$ this gives
$8,15,24$. Golden-ratio-saturated geometries provide examples where all
quadratic invariants remain within $\Z[\varphi]$ by ring closure. The
construction is independent of any specific physical model.

\bibliographystyle{plain}
\bibliography{gentry-shape-space}

\end{document}